\documentclass[12pt]{article}
\usepackage[T1]{fontenc}
\usepackage[utf8]{inputenc}
\usepackage{lmodern}
\usepackage{textcomp}
\usepackage{enumitem}
\usepackage{geometry}
\geometry{margin=1in}
\begin{document}
\begin{center}
Answer Key - Psychology - Undergraduate Level Assessment 4 \
\small (Answers are ordered exactly as in the question paper.)
\end{center}

\section*{Multiple Choice}
\begin{enumerate}[label=\arabic*.]
\item \textbf{Answer:} A
\\ \textit{Options:} A) Researcher; B) therapist; C) faculty member; D) mentor; E) Student
\\ \textit{Note:} Vicarious liability is most likely to be an issue when a psychologist is acting as a researcher because they may be held responsible for the actions of their assistants or participants during the research process, particularly if those actions result in harm or ethical violations.
\item \textbf{Answer:} A
\\ \textit{Options:} A) analytic, spatial, and verbal; B) mathematical, spatial, and analytic; C) verbal, mathematical, and recognizing emotional expressions; D) mathematical, spatial, and musical; E) spatial, musical, and recognizing emotional expressions
\\ \textit{Note:} The left cerebral hemisphere is primarily responsible for analytic thinking, language processing, and verbal communication, which underscores its specialization in these functions.
\item \textbf{Answer:} D
\\ \textit{Options:} A) associative; B) distributive; C) additive; D) disjunctive; E) divisive
\\ \textit{Note:} A disjunctive task is characterized by a group arriving at a solution by choosing the best option provided by any one of its members, rather than requiring all members to contribute equally to the outcome.
\item \textbf{Answer:} A
\\ \textit{Options:} A) Hyperthyroidism; B) Addison's disease; C) Low levels of melatonin; D) An excess of testosterone; E) Hypothyroidism
\\ \textit{Note:} Hyperthyroidism can lead to an overproduction of thyroid hormones, which can cause symptoms such as increased heart rate, nervousness, and irritability that mimic anxiety disorders.
\item \textbf{Answer:} A
\\ \textit{Options:} A) 2.74; B) 4.00; C) 6.08; D) 5.12; E) 3.92
\\ \textit{Note:} The correct answer of 2.74 represents the standard deviation of the sample, which quantifies the amount of variation or dispersion in the data set, calculated using the formula for the sample standard deviation that accounts for the mean and the number of data points.
\end{enumerate}

\section*{True / False}
\begin{enumerate}[label=\arabic*.]
\setcounter{enumi}{5}
\item \textbf{Answer:} False
\\ \textit{Options:} A) True; B) False
\\ \textit{Note:} Fatigue significantly impairs cognitive function, attention, and reaction times, thereby increasing the likelihood of workplace accidents compared to other contributing factors.
\item \textbf{Answer:} False
\\ \textit{Options:} A) True; B) False
\\ \textit{Note:} Allowing a child to remain in the classroom during time-out does not constitute positive punishment, as time-out is intended to remove the child from the reinforcing environment of the classroom to decrease disruptive behavior, rather than increase it.
\item \textbf{Answer:} False
\\ \textit{Options:} A) True; B) False
\\ \textit{Note:} The statement is false because the Wechsler Intelligence Scale for Children (WISC) is designed so that the majority of the population, about 68 percent, scores within one standard deviation of the mean (which is 100), placing the range of scores between 85 and 115, rather than the broader range of 70 to 130.
\item \textbf{Answer:} False
\\ \textit{Options:} A) True; B) False
\\ \textit{Note:} In a skewed distribution with a left tail, the mean is typically lower than the median and mode because it is influenced by the lower values in the tail, making the statement false.
\item \textbf{Answer:} False
\\ \textit{Options:} A) True; B) False
\\ \textit{Note:} The statement is false because the correct deviations from the mean for the values 4, 4, 6, 7, and 9 are -2, -2, 0, 1, and 3 respectively, not -2, 0, 2, 3, and 4.
\end{enumerate}

\section*{Long Form Response}
\begin{enumerate}[label=\arabic*.]
\setcounter{enumi}{10}
\item \textbf{Answer:} Asking for a candidate's age may be justified as a bona fide occupational requirement (BFOQ) in specific circumstances where age is directly linked to the ability to perform essential job functions. For instance, in roles that involve the supervision of minors, such as in educational or childcare settings, certain age-related legal stipulations may necessitate age verification to ensure compliance with safeguarding laws. Additionally, in industries where physical capabilities decline with age, such as in certain types of labor-intensive jobs, age may be relevant to ensure workplace safety. However, it is crucial that employers approach this issue with caution to avoid age discrimination and ensure that age is not used as a blanket exclusion criterion without legitimate justification. Ultimately, the ethical implications hinge on balancing the need for age as a job requirement against the principles of fairness and equality in hiring practices.
\\ \textit{Note:} correct because it identifies specific contexts, such as the supervision of minors and physical capability requirements, where age can be a legitimate criterion for ensuring compliance with legal standards and job performance.
\item \textbf{Answer:} PET scans, or Positron Emission Tomography scans, are instrumental in debunking the myth that humans only utilize 10 percent of their brains. These scans allow researchers to observe brain activity in real-time by tracking the flow of glucose, which is fundamental for brain function. The results consistently show that virtually all parts of the brain have a role and are active at different times, depending on the tasks being performed. This misconception has implications in psychology, as it may lead to misunderstandings about human potential and cognitive capabilities, potentially impacting educational and therapeutic approaches. By highlighting the extensive functionality of the entire brain, PET scans reinforce the notion that all regions contribute to our behaviors, thoughts, and emotions, countering the oversimplified view of brain usage.
\\ \textit{Note:} correct because PET scans demonstrate that all areas of the brain are active during various tasks, thereby disproving the 10 percent myth and highlighting the importance of understanding brain function in psychology.
\end{enumerate}

\vfill
\noindent\textit{End of Answer Key}
\end{document}